\documentclass[../../main.tex]{subfiles}
\begin{document}

\begin{center}
  {\fontsize{48pt}{58pt}\selectfont \textbf{1981}\normalsize\textbf{年度}}
\end{center}

\begin{table}[htb]
  \centering
  \caption{出題テーマと難易度}
  \begin{tabular}{|c|p{9cm}|c|}
    \hline
    \textbf{問題番号} & \textbf{テーマ} & \textbf{難易度} \\
    \hline
    第1問 & 数列と2進法展開 & \\
    \hline
    第2問 & 3次関数と面積の最大値 & \\
    \hline
    第3問 & 円運動と面積の最大値 & \\
    \hline
    第4問 & 定積分と極限 & \\
    \hline
  \end{tabular}
\end{table}

\begin{table}[htb]
  \centering
  \caption{解答}
  \renewcommand{\arraystretch}{1.4}
  \begin{tabular}{|c|c|p{9cm}|}
    \hline
    \textbf{問題番号} & \textbf{小問} & \textbf{答え} \\
    \hline
    第1問 & - & $\alpha=\dfrac{1}{3}$，$a_{2k-1}=\dfrac{1}{3}(4^{k-1}-1),\ a_{2k}=\dfrac{2}{3}(4^{k-1}-1)$ \\
    \hline
    第2問 & - & $\dfrac{50}{9}$ \\
    \hline
    第3問 & - & $\dfrac{\sqrt{409+88\sqrt{22}}}{6}$ \\
    \hline
    第4問 & (1) & $\dfrac{\pi}{2}$ \\
    \cline{2-3}
    & (2) & $0<t\le\dfrac{1}{2}$ または $t=1$ \\
    \hline
  \end{tabular}
\end{table}

\paragraph{試験形態}
4問．試験時間は要確認．

\end{document}
